% !TEX root = ../notes_missingsemester.tex
% Block 10: AI Dev & Vibe Coding
%%%%%%%%%%%%%%%%%%%%%%%%%%%%%%%%%%%%%%%%%%%%%%%%%%%%%%%%%%%%%%%%%%%%%%

% In a first draft make students get a small LLM running locally and set up as a coding assistant.
% Setting this up with an .sh file is fine, but students should get the chance to 
% go the way by foot as well. 




\index{AI-assisted development}\index{vibe coding}\index{LLM}\index{agent}\index{Aider}\index{Qwen}\index{Ollama}
\chapter{AI Dev \& Vibe Coding}
\printchapterversion{10_ai_dev_vibecoding}
\newthought{The tools changed. The fundamentals did not. Know both.}

\marginnote{\newthought{Recommended viewing.} Erik Schluntz on ``Vibe Coding in
Prod''~\cite{Schluntz2025VibeCoding} walks through merging 22{,}000 lines mostly written by Claude
into a production codebase, and the safety workflow that made it tenable. The Claude team's ``How
Anthropic teams use Claude Code''~\cite{AnthropicTeams2025ClaudeCode} surveys how data scientists,
product engineers, lawyers, and marketers at Anthropic actually drive the agent.}

% ==============================================================================================
% BLOCK 1: INTRODUCTION
% ==============================================================================================

\section{Two Ways AI Helps You Code}

\newthought{AI coding tools} come in two flavours, and knowing when to use which matters.

\begin{description}
    \item[Terminal CLI agents]\index{CLI agent} Aider\index{Aider} at \url{https://aider.chat/},
    Claude Code\index{Claude Code} at \url{https://docs.anthropic.com/en/docs/claude-code},
    and Codex CLI work directly in your repository from the terminal. They read your files, propose
    edits, and run commands. Best for implementing features, debugging, and refactoring across
    multiple files.
    \item[IDE-integrated agents]\index{IDE agent} Cursor\index{Cursor} at \url{https://www.cursor.com/},
    Windsurf at \url{https://windsurf.com/}, and Cline at \url{https://cline.bot/} are full editors
    with AI built in, visual diffs, and inline suggestions. Best for editing in context, reviewing
    changes visually, and inline completions while typing.
\end{description}

\marginnote{Autocomplete tools such as GitHub Copilot\index{GitHub Copilot} live at
\url{https://github.com/features/copilot}. Chat tools such as ChatGPT at \url{https://chatgpt.com/}
and Claude at \url{https://claude.ai/} answer questions without touching your files. CLI and IDE
agents are the next level: They read your whole project, not just the current file.}

\newthought{This course uses Aider} as the agent, talking to a self-hosted Qwen2.5-Coder model via
Ollama\index{Ollama} at \url{https://ollama.com/}. The Ollama instance is provided by the course,
no external API account is required, and your prompts never leave the lab network. Qwen2.5-Coder is
a small open-weights coder model. It is fast, free, and good enough to teach the workflow, but it is
not as capable as a frontier hosted model. Everything you learn here transfers to the larger tools,
the constraints just bite a little harder when the model is small.

% ==============================================================================================
\newpage
\section{Optional Exercise: Catch Up to \texttt{v0.8}}\index{exercises}\index{Git!tag}\index{recovery}

\marginnote{\newthought{Tag map.}\index{Git!tag}
\texttt{v0.8} marks the end of Block~9.
This is the final block, no further tag is produced.
Browse the full set at
\url{https://github.com/schutera/pixelwise/tags}.}

\newthought{Catching up via tag.}\index{Git!tag}\index{recovery}
This catch-up extends the one in Block~9. If you have not done that one,
do it first: It lands you at \texttt{v0.7}, the end-of-Block~8 state.
The exercise here picks up from there and rewinds to \texttt{v0.8}, the
end-of-Block~9 state, with structured JSON logging, an upgraded
\texttt{/health} endpoint, a Discord webhook for error-level alerts,
and the \texttt{/api/analytics} endpoint serving aggregates from the
PostgreSQL \texttt{predictions} table.

\newthought{Both VMs this time.} Block~9 added observability and
analytics on top of the running service, so this catch-up brings both
machines up to \texttt{v0.8}. The Git-tracked changes since \texttt{v0.7}
include structured JSON logging in \texttt{app/main.py} that emits
timestamps, levels, and per-request fields, an upgraded \texttt{/health}
endpoint that distinguishes liveness from readiness,
\texttt{scripts/health\_check.sh} for the cron-driven probe, the Discord
webhook integration that posts on error-level logs, the
\texttt{/api/analytics} endpoint reading aggregates from the
\texttt{predictions} table, and a refreshed \texttt{.env.example} with
\texttt{DISCORD\_WEBHOOK\_URL}. The local \texttt{.env} on each VM needs
the new webhook URL written by hand; \texttt{setup-server.sh} then
reconciles dependencies and restarts the service.

\newthought{Run the catch-up on \texttt{dev}.}
Rewind the repository to \texttt{v0.8}, refresh \texttt{.env} from the
updated example and edit it to set \texttt{DISCORD\_WEBHOOK\_URL},
reinstall pinned dependencies if your virtual environment was wiped,
then run \texttt{setup-server.sh}. The exact commands are in the margin.

\marginnote{\newthought{Sequence on \texttt{dev}.}\\[0.4em]
\texttt{cd \textasciitilde/pixelwise}\\
\texttt{git fetch origin -{}-tags}\\
\texttt{git reset -{}-hard v0.8}\\
\texttt{cp .env.example .env}\\
\texttt{nano .env}\\
\texttt{source .venv/bin/activate}\\
\texttt{pip install -r requirements.txt}\\
\texttt{bash setup-server.sh}\\[0.3em]
Edit \texttt{.env} to set \texttt{DISCORD\_WEBHOOK\_URL} and any local
overrides. If \texttt{.venv/} is missing, run
\texttt{python3 -m venv .venv} before the
\texttt{source} step.}

\newthought{Run the catch-up on \texttt{prod}.}
Same sequence on prod, where \texttt{setup-server.sh} also restarts the
service so the new logging, health, and analytics code goes live. After
the script finishes, verify the API is reachable from \texttt{dev}.

\marginnote{\newthought{Sequence on \texttt{prod}.}\\[0.4em]
\texttt{ssh produser@192.168.56.11}\\
\texttt{cd /opt/pixelwise}\\
\texttt{git fetch origin -{}-tags}\\
\texttt{git reset -{}-hard v0.8}\\
\texttt{cp .env.example .env}\\
\texttt{nano .env}\\
\texttt{source .venv/bin/activate}\\
\texttt{pip install -r requirements.txt}\\
\texttt{bash setup-server.sh}\\[0.3em]
Then from \texttt{dev}:
\texttt{curl http://192.168.56.11:8000/health}
should return \texttt{\{"status":"ok"\dots\}}.}

\marginnote{\newthought{Already have a fork?}\index{Git!fork}
If your own fork carries the \texttt{v0.8} tag, swap \texttt{origin} on
\texttt{dev} with
\texttt{git remote set-url origin git@github.com:<you>/pixelwise.git}
so push and pull continue to target your account.
The HTTPS clone on \texttt{prod} stays read-only.}

\newthought{You are now} at the state where Block~9 ended, with both
machines aligned, ready to start Block~10.

% ==============================================================================================
% BLOCK 2: CONCEPTS AND PRACTICE
% ==============================================================================================

\newpage
\section{The Agent Loop}\index{agent loop}

\newthought{How these tools actually work:} The agent loop repeats until the task is done.

\begin{itemize}
    \item Observe: The agent collects context, your repository structure, Git status, recent commits,
    and the files you point it at.
    \item Inspect: It analyses the code and your request to understand what needs to change.
    \item Choose: It selects an action, edit a file, run a command, or ask you a question.
    \item Act: It executes the action, gets the result, and loops back to observe.
\end{itemize}

\newthought{The harness} is the program wrapping the Large Language Model, LLM\index{LLM}, and it
validates every action before execution. Is the tool recognised? Are the arguments valid? Is the
file path inside the repository? The LLM proposes, the harness gatekeeps. In Aider the harness is
deliberately thin: It accepts edit blocks the model emits, applies them with a diff, optionally
runs a test or lint command, and falls back to asking you if anything is unclear.

\newthought{Quality depends entirely on the context you give the agent.} Garbage in, garbage out.
A vague prompt produces vague code. A precise prompt with file references produces precise changes.
With a small local model this matters more, not less: Qwen2.5-Coder will not paper over a sloppy
prompt the way a frontier model sometimes can.

\subsection{Sandboxing}\index{sandboxing}

\newthought{AI agents run on your machine,} so what stops them from doing damage? An unrestricted
agent can run \texttt{rm -rf /}, \texttt{curl} your secrets to an external server, install packages,
or modify system files. It runs with your user's permissions.

\newthought{How sandboxing works in practice:}

\begin{description}
    \item[Code-only by default] Aider edits files but does not run shell commands on its own. If you
    want it to run a test or a script, you ask explicitly with \texttt{/run}, and Aider shows you the
    command before executing.
    \item[Permission prompts] Agents like Claude Code ask before executing each command, and you
    approve or deny. The human stays in the loop.
    \item[Restricted execution] Docker containers, VMs, or chroot jails, the agent runs in an
    isolated environment where it physically cannot access the host system. Your dev VM already
    plays this role for this course.
    \item[Allow and deny lists] Configure which commands or directories the agent can access. Some
    tools let you whitelist specific tools, for example \texttt{pytest} but nothing else.
\end{description}

\marginnote{The principle is least privilege: Give the agent the minimum access it needs. Read your
code, yes. Edit files, yes. Run arbitrary commands, only with explicit approval. Access
\texttt{.env} or SSH keys, never.}

% ==============================================================================================
\newpage
\section{Exercise: Set up Aider}\index{exercises}

\newthought{Set up Aider} on your dev VM. Aider\index{Aider} at \url{https://aider.chat/} is an
open-source CLI agent. The course hosts a Qwen2.5-Coder\index{Qwen} instance via
Ollama\index{Ollama} on the lab network, so you do not need an external API key.

\begin{verbatim}
# Install Aider into your project venv
source .venv/bin/activate
pip install aider-chat

# Point Aider at the course-hosted Ollama
export OLLAMA_API_BASE=http://<course-ollama-host>:11434

# Start Aider against the Qwen2.5-Coder model
cd ~/pixelwise
aider --model ollama/qwen2.5-coder
\end{verbatim}

\marginnote{The exact \texttt{OLLAMA\_API\_BASE} value is posted on \url{material.schutera.com}
together with the model tag in use. If your laptop is on the lab Wi-Fi but Aider hangs on the first
prompt, check the URL and that the Ollama port is reachable with
\texttt{curl \$OLLAMA\_API\_BASE/api/tags}.}

\newthought{Once Aider runs}, you have a working instance of the agent loop sitting at your prompt.
Try \texttt{/help} to see the built-in commands, and \texttt{/add app/main.py} to put a file in
context. The next exercises will use the same session as their starting point.

% ==============================================================================================
\newpage
\section{Exercise: Host Your Own Qwen}\index{exercises}\index{Ollama}\index{Qwen}

\newthought{Now run the model yourself.} The course server is a convenience; owning the model host
means no shared rate limits, no dependency on the lab network, and a setup you can reproduce on
any machine you will use after this course. The trade-off is hardware: A Large Language Model,
LLM, is a memory-hungry process, and a small machine forces a small model.

\newthought{Install Ollama} on the machine that will host the model, either your dev VM or your
laptop. The dev VM is convenient, but a laptop with more Random Access Memory, RAM, or a dedicated
Graphics Processing Unit, GPU, will be noticeably faster.

\begin{verbatim}
# Ubuntu / dev VM
curl -fsSL https://ollama.com/install.sh | sh
sudo systemctl status ollama
\end{verbatim}

\marginnote{macOS and Windows have native installers at \url{https://ollama.com/download}. Ollama
exposes its HTTP API on \url{http://localhost:11434} by default. If you host on your laptop and
run Aider on the dev VM, bind Ollama to a reachable interface with \texttt{OLLAMA\_HOST=0.0.0.0}
and open the port in your firewall.}

\newthought{Pick a model size that fits your hardware.} Qwen2.5-Coder ships at several sizes. Pick
the largest one that comfortably fits in free memory, with headroom for the editor and the
PixelWise service running alongside.

\begin{description}
    \item[\texttt{qwen2.5-coder:0.5b}] About 400 MB on disk, roughly 1 GB of RAM in use. Good
    enough to demonstrate the workflow, frustrating for real edits. Use this only if nothing
    larger fits.
    \item[\texttt{qwen2.5-coder:1.5b}] About 1 GB on disk, roughly 2 GB of RAM in use. Usable for
    single-file edits and short conversations. A sensible floor for the dev VM.
    \item[\texttt{qwen2.5-coder:3b}] About 2 GB on disk, roughly 4 GB of RAM in use. Decent quality
    on small files, slow without a GPU.
    \item[\texttt{qwen2.5-coder:7b}] About 4.7 GB on disk, roughly 8 GB of RAM in use. Matches the
    course-hosted model. The sweet spot if your laptop can carry it.
    \item[\texttt{qwen2.5-coder:14b}] About 9 GB on disk, roughly 16 GB of RAM in use, and a GPU
    helps a lot. Comfortable on multi-file tasks.
\end{description}

\marginnote{Check available memory with \texttt{free -h} on Linux or Activity Monitor on macOS.
Leave headroom: A loaded model plus your editor plus the FastAPI service is more weight than the
model size alone suggests. If the system starts swapping, the model is too big and you should
drop one tier.}

\newthought{Pull and smoke-test the model:}

\begin{verbatim}
ollama pull qwen2.5-coder:1.5b
ollama run qwen2.5-coder:1.5b "write a Python function that reverses a string"
\end{verbatim}

The first \texttt{run} after a fresh \texttt{pull} loads the weights into memory and may take a
minute. Subsequent calls are fast.

\newthought{Point Aider at your local Ollama} by overriding the API base before launching, then
selecting the model tag you just pulled:

\begin{verbatim}
export OLLAMA_API_BASE=http://localhost:11434
aider --model ollama/qwen2.5-coder:1.5b
\end{verbatim}

\newthought{Persist the choice} so future sessions pick it up automatically. Update
\texttt{.aider.conf.yml} in the repository root:

\begin{verbatim}
# .aider.conf.yml
model: ollama/qwen2.5-coder:1.5b
read:
  - CONVENTIONS.md
auto-commits: true
\end{verbatim}

\newthought{Verify} by repeating one prompt from the previous exercise against your self-hosted
model. The answer quality will likely drop compared with the course-hosted 7B model, and that is
the point. You now feel the difference between model sizes in your hands, and you know how to
scale up when a task demands it.

% ==============================================================================================
\newpage
\section{Git as a Safety Net}\index{Git!safety net}

\newthought{Commit before you let an agent touch your code.} This is not just an AI workflow, it is
how you work with any risky change: Experimental refactors, dependency upgrades, merge conflict
resolution.

\begin{verbatim}
# Save current state
git add -A && git commit -m "pre-agent checkpoint"

# See what the agent changed
git diff

# Revert everything if the agent made a mess
git checkout .

# Stash your current state before experimenting
git stash

# Let the agent work on a branch
git checkout -b experiment
# ... agent does its work ...
# If it works, merge. If not, delete the branch.
git checkout main
git branch -D experiment
\end{verbatim}

\marginnote{Always commit clean state first, then change, then review the diff. This pattern
predates AI tools. It is how experienced developers have always worked with risky changes. Aider
will, by default, also auto-commit each successful edit on its own branch of history, so you end
up with a fine-grained log to bisect when something goes wrong.}

\subsection{Reading Diffs}\index{diff}\index{code review}

\newthought{The core review skill.} When an agent or a colleague makes changes, you do not re-read
the entire file. You read the diff: What was added, what was removed, what was modified.

\begin{verbatim}
git diff                    # in the terminal
\end{verbatim}

\newthought{\texttt{git diff} in the terminal,} side-by-side diffs in the IDE, and GitHub's pull
request diff view all exercise the same skill. Students need to read diffs fluently. This is how
code review actually works.

\marginnote{Lines starting with \texttt{+} were added. Lines starting with \texttt{-} were removed.
Lines starting with a space are unchanged context. The \texttt{@@} header tells you which file and
which line numbers changed.}

% ==============================================================================================
\newpage
\section{Exercise: Commit Clean State}\index{exercises}

\newthought{First rule, commit clean state} before letting the agent touch anything:

\begin{verbatim}
git add -A
git commit -m "pre-agent checkpoint"
\end{verbatim}

This is your undo button for everything that follows. From this clean state, every change the
agent makes shows up cleanly in \texttt{git diff} and can be reverted with a single command.

\newthought{Verify the safety net} works. Let the agent make a small, deliberately throwaway
change, then read the diff and revert:

\begin{verbatim}
git diff
git checkout .
\end{verbatim}

You should land back at the pre-agent checkpoint with the agent's edits gone. Until you can do
this without thinking, do not give the agent anything important to touch.

\marginnote{If Aider's auto-commit is on, \texttt{git checkout .} alone will not undo committed
edits. Use \texttt{git reset -{}-hard <pre-agent-sha>} to rewind, or pass
\texttt{-{}-no-auto-commits} when you launch Aider so the safety net behaves the way the manual
checkpoint workflow assumes.}

% ==============================================================================================
\newpage
\section{Context Management}\index{context window}\index{context management}

\newthought{Agents have a limited context window,} which is how much text the LLM can process at
once. Qwen2.5-Coder running on the course server has a smaller window than a frontier hosted model,
so the budget bites earlier. Long conversations degrade quality: The agent forgets earlier
instructions, repeats mistakes, or contradicts itself. Older parts of the conversation get
compressed or dropped.

\newthought{Practical rules:}
\begin{itemize}
    \item Start a fresh session when switching tasks. \texttt{/clear} drops the chat history,
    \texttt{/reset} drops both the history and the files in context.
    \item Point the agent at specific files rather than letting it read everything. In Aider:
    \texttt{/add app/main.py} to put a file in scope for edits, \texttt{/read README.md} for files
    you want it to know about but not modify, and \texttt{/drop <file>} to evict files you no longer
    need.
    \item Keep instructions files short and focused. The cost of every token in your conventions
    file is paid on every turn.
\end{itemize}

\marginnote{A small model is a forcing function: It rewards tight scoping and punishes context
bloat almost immediately. The habits you build here scale up to larger models, but the larger
models make the habits optional. Build them now.}

% ==============================================================================================
\newpage
\section{Exercise: Explore the Tool}\index{exercises}

\newthought{Get comfortable with the conversation loop} by asking Aider to explain a file, describe
the PixelWise architecture, or find where a function is called. Use \texttt{/ask} mode for
read-only conversation so the agent does not try to edit anything while you are exploring. Try:

\begin{itemize}
    \item \texttt{/add app/classifier.py} then \texttt{/ask explain what this file does}
    \item \texttt{/ask where is classify\_batch called?}
    \item \texttt{/ask describe the full request flow from browser to database}
\end{itemize}

\newthought{Watch the context fill up.} Notice how the agent's answers degrade if you keep piling
questions into the same conversation, and how a fresh \texttt{/clear}'d session with a single,
file-scoped prompt produces a sharper answer. Pointing the agent at specific files rather than
asking it to look at the whole project is the cheapest quality gain you have, and it is the
difference between a useful and a useless small model.

% ==============================================================================================
\newpage
\section{Teaching the AI Your Project}\index{instructions file}\index{CONVENTIONS.md}

\newthought{Instructions as code:} Declarative files that persist across sessions. You write them
once, the AI follows them every time. This is the highest-leverage thing you can do with AI tools,
and it is even more important with a small model than with a large one, because a small model has
less general knowledge to fall back on when the project diverges from defaults.

\begin{description}
    \item[\texttt{CONVENTIONS.md}]\index{CONVENTIONS.md} Project-level instructions Aider auto-reads
    when launched with \texttt{aider -{}-read CONVENTIONS.md}, or whenever the path is listed under
    \texttt{read:} in \texttt{.aider.conf.yml}. Cover conventions, architecture, and what NOT to do.
    \item[\texttt{.aider.conf.yml}]\index{.aider.conf.yml} Aider's per-project config. Use it to pin
    the model, declare auto-loaded read-only files, set \texttt{auto-commits}, and lock in the lint
    and test commands.
    \item[\texttt{CLAUDE.md} and \texttt{.cursorrules}] The equivalent files for Claude Code and
    Cursor. Different filename, identical idea: A markdown file in your repository that the AI
    reads before starting work.
\end{description}

\newthought{The pattern across all tools} is the same: A short, declarative, human-readable file
that tells the agent what your repository is and how to behave inside it. The agent loads it once,
applies it on every turn.

\subsection{Task Playbooks}\index{playbook}

\newthought{Playbooks are reusable recipes} for specific tasks: Structured \texttt{.md} files with
step-by-step instructions the AI follows consistently. Example: How to add a new API endpoint to
this project. Store them under \texttt{docs/playbooks/} and pull one into context with
\texttt{/read docs/playbooks/add-endpoint.md} at the start of the relevant session. The agent
should then produce correct code from the playbook and the task description alone.

\marginnote{Larger tools have built-in versions of this idea, for example Claude Code's skill files.
With Aider and a local model, plain markdown files plus \texttt{/read} cover the same need without
extra machinery.}

% ==============================================================================================
\newpage
\section{Exercise: Write a CONVENTIONS.md}\index{exercises}

\newthought{Write a \texttt{CONVENTIONS.md}} for the PixelWise project. Document the stack, the
conventions, and what the AI should and should not do. Example structure:

\begin{verbatim}
# PixelWise

## Stack
- Python 3.11, FastAPI, SQLAlchemy, PostgreSQL
- Nginx reverse proxy, vanilla JS frontend
- sklearn LogisticRegression model on MNIST

## Conventions
- All config via .env, never hardcode secrets
- Batch-first inference, classify_batch is the core
- Pydantic models for all API contracts

## Do NOT
- Modify .env or expose secrets
- Add new dependencies without updating requirements.txt
- Change the classifier interface without updating tests
\end{verbatim}

Commit it to the repository, then wire Aider up so it loads automatically:

\begin{verbatim}
# .aider.conf.yml
model: ollama/qwen2.5-coder
read:
  - CONVENTIONS.md
auto-commits: true
\end{verbatim}

From now on, every Aider session in this project starts with \texttt{CONVENTIONS.md} already in
context.

% ==============================================================================================
\newpage
\section{Exercise: Write a Task Playbook}\index{exercises}

\newthought{Create a \texttt{docs/playbooks/} directory} and write a playbook for one specific
task, for example how to add a new API endpoint:

\begin{verbatim}
mkdir -p docs/playbooks
nano docs/playbooks/add-endpoint.md
\end{verbatim}

Include which files to modify, which patterns to follow, and which tests to run afterwards. Then
test it from a clean Aider session:

\begin{verbatim}
aider
> /read docs/playbooks/add-endpoint.md
> add a GET /api/status endpoint that returns build version and uptime
\end{verbatim}

\newthought{Iterate on the playbook} until the agent can carry out the task from the playbook alone,
without you correcting it mid-way. A good playbook is one a tired colleague could also follow.

% ==============================================================================================
\newpage
\section{Working with AI Agents Effectively}

\newthought{The agent-edit-verify loop:}
\begin{itemize}
    \item Agent makes changes.
    \item \texttt{git diff}, review every line.
    \item Run the app, it works or it does not.
    \item If not, feed the error back to the agent or fix it yourself.
    \item Iterate until it is right, or revert and try a different approach.
\end{itemize}

\newthought{Use \texttt{/architect}} to let Aider plan before it edits. In architect mode the model
proposes a structured plan and lists the files it would touch, but does not write the changes until
you accept. Review the plan, adjust, then let it execute. Do not jump straight to implementation,
especially with a small model that can otherwise commit to the wrong design early.

\newthought{Code review with AI:} Have the agent review your own code. It catches things you are
blind to after staring at the same file for hours.

% ==============================================================================================
\newpage
\section{Exercise: Make a Real Change with the AI}\index{exercises}

\newthought{Student's choice:} A new feature, a refactor, or a bug fix. Follow the
agent-edit-verify loop:

\begin{itemize}
    \item Use \texttt{/architect} first to scope the change before any edits.
    \item Let the agent make changes.
    \item Review with \texttt{git diff}, read every line.
    \item Run the app, does it work?
    \item If not, feed the error back. Iterate or revert.
\end{itemize}

\newthought{Stay in the loop} the entire time. The moment you stop reading the diffs is the moment
the agent starts shipping bugs into your repository with your name on the commit.

% ==============================================================================================
\newpage
\section{Exercise: Branching Workflow}\index{exercises}

\newthought{Create a branch for the AI's work} so the main line stays clean while you experiment:

\begin{verbatim}
git checkout -b ai-experiment
# ... let the agent work ...
git diff main
# If good: merge. If bad: delete the branch.
\end{verbatim}

The branch is the boundary between ``the agent has free rein'' and ``the codebase you trust''.
Merge only what passes the diff and the run.

% ==============================================================================================
\newpage
\section{Exercise: Group Discussion}\index{exercises}

\newthought{Compare notes with the room.} When was Aider faster than writing code by hand? When was
it slower? When was it confidently wrong? What did the \texttt{CONVENTIONS.md} change about the
agent's output quality? Where did the small model's context limit show up first?

\newthought{Name one thing} you will keep doing, one thing you will stop doing, and one habit you
want to acquire before the agent touches production code on your behalf.

% ==============================================================================================
\newpage
\section{Optional: MCP and Tool Use}\index{MCP}\index{Model Context Protocol}

\newthought{Agents can do more than edit code.} The Model Context Protocol, MCP\index{MCP}, lets an
agent connect to external tools: Reading documentation, searching the web, running tests, querying
databases, interacting with APIs. The agent is no longer limited to what is in your repository.

\marginnote{MCP reference: \url{https://modelcontextprotocol.io/}. Aider with a local Qwen model
does not ship with MCP integrations, so this is forward-looking context rather than something to
wire up in this course. The frontier hosted agents are where this currently lives.}

% ==============================================================================================
\newpage
\section{Optional: When to Trust AI Output}\index{AI trust}

\begin{description}
    \item[Strong] Boilerplate, tests for known behaviour, documentation, refactoring, and
    explaining unfamiliar code.
    \item[Weak] Security-sensitive code, novel algorithms, anything touching infrastructure or
    secrets.
    \item[Never] Blind trust without reading the output. With a small local model the floor is
    lower, so even on the strong cases you read every line.
\end{description}

% ==============================================================================================
\newpage
\section{Security Thread}

\newthought{AI-generated code must pass the same security review as human-written code.}
Prompt injection\index{prompt injection} is the trap: User-submitted text that ends up in an AI
prompt is an injection vector, never pipe untrusted input into an AI call. The AI must never be
given access to \texttt{.env}, SSH keys, or production credentials directly. Aider's
\texttt{/read} command is for files you want the model to see; treat anything in \texttt{.env} or
\texttt{\textasciitilde/.ssh/} as off-limits to that command.

\newthought{A note on what comes next.} This is the final block. The course wrap-up sits in the
recap below: Every layer you touched, from two empty VMs to a deployed, monitored, AI-assisted
service, is now yours to maintain.

% ==============================================================================================
% BLOCK 3: CONSOLIDATION
% ==============================================================================================

\section{Self-Reflection and Recap}

\newthought{Reflection questions:}
\begin{itemize}
    \item Why do you commit clean state before letting an agent touch your code?
    \item What is the difference between a CLI agent and an IDE agent? When would you use each?
    \item How does a \texttt{CONVENTIONS.md} improve the quality of AI-generated code?
    \item Where did the small local model break down first, context, reasoning, or knowledge of
    your stack?
    \item What are the risks of giving an AI agent unrestricted access to your machine?
    \item When should you trust AI output, and when should you be sceptical?
\end{itemize}

\newthought{Recap} of key concepts:
\begin{itemize}
    \item The agent loop of observe, inspect, choose, act, and the harness that gatekeeps every
    action.
    \item Git as the safety net: Clean checkpoint, branch, diff, revert.
    \item Sandboxing and least privilege: The agent gets the minimum access it needs.
    \item Context management: Short conversations, file-scoped prompts, \texttt{/add},
    \texttt{/read}, \texttt{/drop}, \texttt{/clear}, instructions files.
    \item \texttt{CONVENTIONS.md} and task playbooks as the highest-leverage way to steer AI output.
\end{itemize}

\marginnote{\newthought{Milestone.} Students have used a terminal AI agent backed by a self-hosted
small model on a real codebase. \texttt{CONVENTIONS.md} and a task playbook are committed to the
repository. Every student can read a diff, revert a bad change, and iterate with an agent.}

\marginnote{\newthought{Course wrap-up.} Look at what you have built, from two empty VMs to a
deployed, monitored, auto-deploying AI web service with structured logging, database persistence,
CI/CD, and now AI-assisted development. The full stack, end to end. You own every layer.}
